\chapter{Les chromosomes}\label{ch:g9:chromosomes}

\Cref{rem:g9:heredity-and-traits:question} a acculé les
déterminants dans le \omterm{def:g6:cell-unit-of-life:parts}{noyau}. Attrape une \omterm{def:g5:first-look-at-cells:cell}{cellule} au bon
moment et le \omterm{def:g6:cell-unit-of-life:parts}{noyau} abat son jeu~: son contenu s'enroule en
corps dénombrables et photographiables --- les \omterm{def:g9:chromosomes:chromosomes}{chromosomes}.
Les compter, dans les \omterm{def:g5:first-look-at-cells:cell}{cellules} du corps et dans les
\omterm{def:g8:sexual-reproduction:gametes}{gamètes}, explique presque tout ce que réclamait l'album de
famille~: les moitiés, les paires, le brassage, et le pile
ou face qui décide si un bébé sera un garçon ou une fille.

\section{Le noyau, pris sur le fait}

\begin{definition}[Les chromosomes]\label{def:g9:chromosomes:chromosomes}
Un \emph{chromosome}\index{chromosome} est un corps en forme
de fil, dans le \omterm{def:g6:cell-unit-of-life:parts}{noyau}, visible (bien coloré, au moment de
la division) quand le contenu du \omterm{def:g6:cell-unit-of-life:parts}{noyau} s'enroule pour le
transport. Entre deux divisions, les mêmes fils sont
déroulés et invisibles --- état de travail~; à la division,
ils se condensent --- état de déplacement. Photographiés et
rangés, les chromosomes d'une \omterm{def:g5:first-look-at-cells:cell}{cellule} forment son
\emph{caryotype}\index{caryotype}~: le jeu complet, étalé.
\end{definition}

\begin{proposition}[Les comptes humains]\label{prop:g9:chromosomes:counts}
Les \omterm{def:g9:chromosomes:chromosomes}{caryotypes} livrent trois nombres qui portent tout le
chapitre~:
\begin{enumerate}
  \item toute \omterm{def:g5:first-look-at-cells:cell}{cellule} du corps humain contient \emph{46}
        \omterm{def:g9:chromosomes:chromosomes}{chromosomes} --- \omterm{def:g1:the-five-senses:sense}{peau}, \omterm{prop:g7:digestion-and-nutrients:absorption}{foie}, \omterm{def:g8:nervous-communication:neuron}{neurone} pareillement
        (les divisions de
        \cref{prop:g5:first-look-at-cells:growth} copient
        fidèlement le jeu dans chaque \omterm{def:g5:first-look-at-cells:cell}{cellule})~;
  \item une fois étalés, les 46 se rangent en
        \emph{23 paires}, les deux partenaires d'une paire
        s'accordant par la taille et les bandes~;
  \item les \emph{\omterm{def:g8:sexual-reproduction:gametes}{gamètes}} font exception~: la \omterm{def:g4:animal-reproduction:sexual}{cellule œuf}
        et la \omterm{def:g8:reproductive-systems:male}{cellule mâle} en portent \emph{23} --- un
        \omterm{def:g9:chromosomes:chromosomes}{chromosome} de chaque paire, aucune paire complète.
\end{enumerate}
\end{proposition}
\admitted

\begin{example}[L'arithmétique de la fécondation]\label{ex:g9:chromosomes:arithmetic}
Les comptes s'emboîtent comme une démonstration. Les
\omterm{def:g8:sexual-reproduction:gametes}{gamètes} à 23 chacun~; la \omterm{def:g5:flower-to-fruit:steps}{fécondation}
(\cref{prop:g8:sexual-reproduction:core}) les additionne~:
la \omterm{def:g4:animal-reproduction:sexual}{cellule œuf} fécondée en contient $23 + 23 = 46$ --- les
paires rétablies, un partenaire de chaque paire venant de
la mère, l'autre du père. Chaque \omterm{def:g5:first-look-at-cells:cell}{cellule} du corps de
l'enfant porte alors l'apport des deux parents --- ce qui
est exactement ce que \cref{rem:g4:animal-reproduction:mix}
observait de l'extérieur, le poulain et sa liste, il y a une
génération de chapitres.
\end{example}

\begin{remark}[Pourquoi les gamètes doivent réduire de moitié]\label{rem:g9:chromosomes:whyhalve}
Supposons que non~: des \omterm{def:g8:sexual-reproduction:gametes}{gamètes} à 46 \omterm{def:g9:chromosomes:chromosomes}{chromosomes} feraient
des enfants à 92, des petits-enfants à 184 --- le jeu
doublant à chaque génération. La division réductrice
particulière des \omterm{def:g8:sexual-reproduction:gametes}{gamètes} --- chacun recevant \emph{un
partenaire de chaque paire} --- est ce qui maintient le
compte de l'\omterm{def:g5:biodiversity:species}{espèce} à 46, génération après génération. Et
remarque la prime~: \emph{quel} partenaire de chaque paire
reçoit un \omterm{def:g8:sexual-reproduction:gametes}{gamète} se règle indépendamment, paire par paire
--- $2^{23}$ jeux possibles à partir d'un seul parent, le
brassage de \cref{prop:g8:sexual-reproduction:shuffle}
enfin compté.
\end{remark}

\section{Les chromosomes portent les déterminants}

\begin{proposition}[Les preuves que le jeu compte]\label{prop:g9:chromosomes:evidence}
Que les \omterm{def:g9:chromosomes:chromosomes}{chromosomes} \emph{soient} les véhicules des
déterminants n'est pas une supposition, mais une conclusion
de preuves~:
\begin{itemize}
  \item ils se comportent exactement comme l'exige
        \cref{prop:g9:heredity-and-traits:conclusions}~:
        présents dans chaque \omterm{def:g5:first-look-at-cells:cell}{cellule}, réduits de moitié
        dans les \omterm{def:g8:sexual-reproduction:gametes}{gamètes}, rétablis à la \omterm{def:g5:flower-to-fruit:steps}{fécondation},
        sélectionnables partenaire par partenaire~;
  \item les accidents de \omterm{def:g9:chromosomes:chromosomes}{caryotype} prouvent la cargaison~:
        un jeu avec un \omterm{def:g9:chromosomes:chromosomes}{chromosome} en trop ou en moins
        change le développement du corps tout entier ---
        un exemplaire supplémentaire de la petite paire 21,
        par exemple, produit l'ensemble reconnaissable de
        \omterm{def:g9:heredity-and-traits:traits}{caractères} appelé trisomie 21. Un seul fil mal
        compté, beaucoup de \omterm{def:g9:heredity-and-traits:traits}{caractères} déplacés~: les fils
        portent bien les consignes des \omterm{def:g9:heredity-and-traits:traits}{caractères}~;
  \item et une paire annonce sa cargaison dans tout
        \omterm{def:g9:chromosomes:chromosomes}{caryotype}~: les \emph{\omterm{def:g9:chromosomes:chromosomes}{chromosomes} sexuels}.
\end{itemize}
\end{proposition}
\admitted

\begin{definition}[X, Y et le pile ou face]\label{def:g9:chromosomes:xy}
La paire 23 existe en deux versions~: deux \omterm{def:g9:chromosomes:chromosomes}{chromosomes}
\emph{X} identiques --- le \omterm{def:g9:chromosomes:chromosomes}{caryotype} d'une fille~; ou un X
et un petit \emph{Y} --- un garçon. La réduction de moitié
joue alors le pile ou face~: toute \omterm{def:g4:animal-reproduction:sexual}{cellule œuf} porte un X~;
une \omterm{def:g8:reproductive-systems:male}{cellule mâle} porte X ou Y, moitié-moitié. Une
\omterm{def:g5:flower-to-fruit:steps}{fécondation} par une \omterm{def:g8:reproductive-systems:male}{cellule mâle} à X fait XX, par une
\omterm{def:g8:reproductive-systems:male}{cellule mâle} à Y fait XY --- une chance sur deux (le
langage du hasard de \cref{ch:g9:heredity-and-traits},
enfin muni de son mécanisme), décidée par le nageur qui
arrive, à l'instant de la \omterm{def:g5:flower-to-fruit:steps}{fécondation}.
\end{definition}

\begin{omfigure}
\begin{tikzpicture}[scale=0.9]
  % mother
  \node[draw, thick, omDef, rounded corners=4pt, font=\small,
    align=center] (mo) at (1.2,4.2) {la mère\\ XX};
  \node[draw, thick, omDef, rounded corners=4pt, font=\small,
    align=center] (fa) at (7.8,4.2) {le père\\ XY};
  % gametes
  \node[draw, omDef, circle, font=\footnotesize] (e1) at (1.2,2.4) {X};
  \node[draw, omDef, circle, font=\footnotesize] (s1) at (6.6,2.4) {X};
  \node[draw, omDef, circle, font=\footnotesize] (s2) at (9.0,2.4) {Y};
  \draw[->, thick] (mo) -- (e1);
  \draw[->, thick] (fa) -- (s1);
  \draw[->, thick] (fa) -- (s2);
  \node[font=\footnotesize, left] at (0.8,2.4) {toute cellule œuf};
  \node[font=\footnotesize, right, align=left] at (9.4,2.4)
    {cellules mâles~:\\ moitié X, moitié Y};
  % offspring
  \node[draw, thick, omThm, rounded corners=4pt, font=\small,
    align=center] (girl) at (3.4,0.4) {XX\\ une fille};
  \node[draw, thick, omMeth, rounded corners=4pt, font=\small,
    align=center] (boy) at (5.8,0.4) {XY\\ un garçon};
  \draw[->, thick, omExo] (e1) -- (girl);
  \draw[->, thick, omExo] (s1) -- (girl);
  \draw[->, thick, omExo] (e1) -- (boy);
  \draw[->, thick, omExo] (s2) -- (boy);
  \node[font=\footnotesize, align=center] at (4.6,-0.5)
    {une chance sur deux, à la fécondation};
\end{tikzpicture}

{\small Le pile ou face de la 23\textsuperscript{e} paire~: le X de la
\omterm{def:g4:animal-reproduction:sexual}{cellule œuf} rencontre une \omterm{def:g8:reproductive-systems:male}{cellule mâle} à X ou à Y --- et le
\omterm{def:g9:chromosomes:chromosomes}{caryotype} est décidé avec la rencontre.}
\end{omfigure}

\begin{example}[Vieilles questions, réponses courtes]\label{ex:g9:chromosomes:answers}
Le jeu règle de vieux comptes. Qu'un enfant soit un garçon
ou une fille ne doit rien aux souhaits, aux régimes ni aux
saisons de l'un ou l'autre parent --- c'est le pile ou
face, joué dans le mécanisme de
\cref{def:g9:chromosomes:xy} (et, pour le registre des
injustices anciennes~: la différence qui décide voyage dans
la \emph{\omterm{def:g8:reproductive-systems:male}{cellule mâle}}). Les moitiés égales répondent aussi
à l'album de famille~: la mère et le père apportent chacun
un demi-jeu complet --- 23 et 23 --- et l'arithmétique de
l'\omterm{def:g9:heredity-and-traits:heredity}{hérédité} honore les deux lignées à égalité.
\end{example}

\begin{method}[Lire un caryotype]\label{met:g9:chromosomes:read}
Devant une photographie de \omterm{def:g9:chromosomes:chromosomes}{caryotype}~:
\begin{enumerate}
  \item compte~: 46 attendus dans une \omterm{def:g5:first-look-at-cells:cell}{cellule} du corps, 23
        dans un \omterm{def:g8:sexual-reproduction:gametes}{gamète}~;
  \item apparie~: les partenaires accordés par la taille et
        les bandes, 22 paires ordinaires plus la 23\textsuperscript{e}~;
  \item lis la 23\textsuperscript{e}~: XX ou XY~;
  \item vérifie les comptes~: toute paire à un ou trois
        éléments explique, d'après
        \cref{prop:g9:chromosomes:evidence}, un
        développement déplacé.
\end{enumerate}
\end{method}

\section{Exercices}

\begin{exercise}[$\star$]\label{exo:g9:chromosomes:1}
Qu'est-ce qu'un \omterm{def:g9:chromosomes:chromosomes}{chromosome}, et quand est-il visible~?
\end{exercise}

\begin{exercise}[$\star$]\label{exo:g9:chromosomes:2}
Donne les trois comptes humains, et dis où chacun vaut.
\end{exercise}

\begin{exercise}[$\star$]\label{exo:g9:chromosomes:3}
Fais l'arithmétique de la \omterm{def:g5:flower-to-fruit:steps}{fécondation}, et dis ce qu'elle
rétablit.
\end{exercise}

\begin{exercise}[$\star$]\label{exo:g9:chromosomes:4}
Pourquoi les \omterm{def:g8:sexual-reproduction:gametes}{gamètes} doivent-ils réduire le jeu de moitié~?
Montre l'emballement sinon.
\end{exercise}

\begin{exercise}[$\star$]\label{exo:g9:chromosomes:5}
Quelles sont les deux versions de la paire 23, et le \omterm{def:g8:sexual-reproduction:gametes}{gamète}
de quel parent tranche entre elles~?
\end{exercise}

\begin{exercise}[$\star$]\label{exo:g9:chromosomes:6}
Que produit un exemplaire supplémentaire du \omterm{def:g9:chromosomes:chromosomes}{chromosome} 21,
et qu'est-ce que cela prouve au sujet des fils~?
\end{exercise}

\begin{exercise}[$\star\star$]\label{exo:g9:chromosomes:7}
Où se produit exactement le brassage de
\cref{prop:g8:sexual-reproduction:shuffle}, dans les termes
de ce chapitre --- et quelle est la taille de l'espace de
brassage d'un seul parent~?
\end{exercise}

\begin{exercise}[$\star\star$]\label{exo:g9:chromosomes:8}
Associe chaque exigence de
\cref{prop:g9:heredity-and-traits:conclusions} au
comportement des \omterm{def:g9:chromosomes:chromosomes}{chromosomes} qui la satisfait.
\end{exercise}

\begin{exercise}[$\star\star$]\label{exo:g9:chromosomes:9}
«~La probabilité d'avoir un garçon est d'une sur deux à
chaque \omterm{def:g2:born-grow-die:cycle}{naissance}, quoi qu'il se soit passé avant.~» Défends
la phrase par le mécanisme --- pourquoi une série de trois
filles ne change-t-elle rien~?
\end{exercise}

\begin{exercise}[$\star\star$]\label{exo:g9:chromosomes:10}
Applique \cref{met:g9:chromosomes:read} à~: (a) une \omterm{def:g5:first-look-at-cells:cell}{cellule}
du corps 46,XX~; (b) un \omterm{def:g8:sexual-reproduction:gametes}{gamète} 23,Y~; (c) une \omterm{def:g5:first-look-at-cells:cell}{cellule} du
corps à 47, l'exemplaire en trop dans la paire 21.
\end{exercise}

\begin{exercise}[$\star\star$]\label{exo:g9:chromosomes:11}
L'histoire a reproché aux reines les filles des royaumes.
Corrige le registre avec \cref{ex:g9:chromosomes:answers},
\omterm{def:g8:sexual-reproduction:gametes}{gamète} par \omterm{def:g8:sexual-reproduction:gametes}{gamète}.
\end{exercise}

\begin{exercise}[$\star\star\star$]\label{exo:g9:chromosomes:12}
Chaque \omterm{def:g5:first-look-at-cells:cell}{cellule} du corps porte les 46 au complet --- et
pourtant un \omterm{def:g8:nervous-communication:neuron}{neurone} et une \omterm{def:g5:first-look-at-cells:cell}{cellule} de \omterm{def:g1:the-five-senses:sense}{peau} diffèrent du tout
au tout (\cref{rem:g6:cell-unit-of-life:twoways}). Formule
l'énigme que cela pose, et la direction de sa résolution
(l'affaire du prochain chapitre~: ce qui, sur un
\omterm{def:g9:chromosomes:chromosomes}{chromosome}, se \emph{lit}).
\end{exercise}

\section{Problème~: la paillasse de cytogénétique}

\begin{problem}[{Problème du week-end --- quatre caryotypes
sur une même table lumineuse}]\label{pb:g9:chromosomes:1}
La table pédagogique d'un laboratoire hospitalier de
cytogénétique~: \omterm{def:g9:chromosomes:chromosomes}{caryotype} A --- 46 \omterm{def:g9:chromosomes:chromosomes}{chromosomes}, 23\textsuperscript{e}
paire XY~; B --- 46, 23\textsuperscript{e} paire XX~; C --- 23
\omterm{def:g9:chromosomes:chromosomes}{chromosomes}, dont un X~; D --- 47, avec trois exemplaires
à la paire 21, 23\textsuperscript{e} paire XX.

\medskip\noindent\textbf{Partie I --- Les lectures.}
\begin{enumerate}
  \item Lis A, B et C avec
        \cref{met:g9:chromosomes:read}~: de qui chacun
        pourrait-il être la \omterm{def:g5:first-look-at-cells:cell}{cellule}~?
  \item Le compte de C n'est pas une erreur. Quelle sorte
        de \omterm{def:g5:first-look-at-cells:cell}{cellule} est-ce, et quelle histoire de division
        son 23 enregistre-t-il~?
  \item Lis D~: le compte, l'appariement, et le
        développement qu'il annonce.
  \item Le \omterm{def:g8:sexual-reproduction:gametes}{gamète} de quel parent a décidé la 23\textsuperscript{e}
        paire de A, et quelles étaient les chances à cette
        \omterm{def:g5:flower-to-fruit:steps}{fécondation}~?
\end{enumerate}

\medskip\noindent\textbf{Partie II --- Les questions
d'arithmétique.}
\begin{enumerate}[resume]
  \item Une stagiaire demande pourquoi les \omterm{def:g5:first-look-at-cells:cell}{cellules} du
        corps d'une même personne s'accordent toutes avec
        le \omterm{def:g9:chromosomes:chromosomes}{caryotype} A. Réponds par les divisions de
        recopiage de
        \cref{prop:g5:first-look-at-cells:growth}.
  \item Une autre demande comment deux parents à 46 ont
        fait le 23 de C. Nomme la division particulière et
        sa règle.
  \item Combine C avec une \omterm{def:g4:animal-reproduction:sexual}{cellule œuf} porteuse de X, puis
        avec sa voisine de laboratoire porteuse de Y ---
        attends~: C \emph{est} sans Y. Énonce ce que le X
        de C garantit sur tout enfant qu'il contribuerait à
        faire.
  \item Calcule les jeux possibles qu'un parent peut
        distribuer~: énonce l'indépendance paire par paire
        et la conclusion en $2^{23}$ --- puis ajoute
        l'appariement, par la \omterm{def:g5:flower-to-fruit:steps}{fécondation}, de deux
        distributions de ce genre.
\end{enumerate}

\medskip\noindent\textbf{Partie III --- La salle de
consultation.}
\begin{enumerate}[resume]
  \item Les parents d'un enfant au \omterm{def:g9:chromosomes:chromosomes}{caryotype} D demandent~:
        «~qu'avons-nous fait de mal~?~» Donne la réponse
        honnête du mécanisme sur les accidents de la
        réduction de moitié --- et dis ce que les 47 fils
        de l'enfant changent et ne changent pas à l'amour
        présent dans la pièce. (Réponds comme la
        conseillère, avec la \omterm{def:g1:living-or-not:living}{biologie} sans détour.)
  \item Un couple avec trois filles demande les chances
        d'avoir un fils ensuite. Réponds par le pile ou
        face, et corrige l'intuition de la série de chance.
  \item Un visiteur affirme qu'un régime peut choisir le
        sexe d'un bébé. Passe l'affirmation au crible de
        \cref{def:g9:chromosomes:xy}~: où la décision
        est-elle prise, et que pourrait atteindre un
        régime~?
  \item Ferme la paillasse en une phrase~: ce que les
        quatre photographies prouvent ensemble sur le
        véhicule de l'\omterm{def:g9:heredity-and-traits:heredity}{hérédité}.
\end{enumerate}
\end{problem}
